\documentclass[11pt,a4paper]{article}

%========================
% PACKAGES
%========================
\usepackage[utf8]{inputenc}
\usepackage[T1]{fontenc}
\usepackage[french]{babel}
\usepackage[a4paper,margin=1.8cm]{geometry}
\usepackage{amsmath,amssymb,amsthm}
\usepackage{lmodern}
\usepackage{xcolor}
\usepackage{tikz}
\usepackage{fancyhdr}

%========================
% HEADER
%========================
\pagestyle{fancy}
\fancyhf{}
\fancyhead[L]{Physique-Chimie 3\ieme}
\fancyhead[C]{Mécanique}
\fancyhead[R]{\thepage}

%========================
% ENVIRONNEMENTS
%========================
\newtheorem{definition}{Définition}
\newtheorem{property}{Propriété}
\newtheorem{example}{Exemple}
\newtheorem{remark}{Remarque}

\title{\Huge Mécanique – Cours complet niveau 3\ieme}
\date{}

\begin{document}

\maketitle

%====================================================
\section{Pourquoi étudier la mécanique ?}

La mécanique permet de décrire et comprendre les mouvements :
\begin{itemize}
\item voitures
\item planètes
\item sportifs
\item objets en chute
\end{itemize}

\textbf{Idée clé :}  
\begin{center}
\textit{Décrire un mouvement + expliquer pourquoi il se produit}
\end{center}

%====================================================
\section{Référentiel et mouvement}

\begin{definition}
Un référentiel est un objet par rapport auquel on décrit un mouvement.
\end{definition}

\begin{example}
\begin{itemize}
\item Dans un train → tu es immobile
\item Pour quelqu’un sur le quai → tu es en mouvement
\end{itemize}
\end{example}

\textbf{Conclusion :}  
Le mouvement dépend du point de vue.

%====================================================
\section{Trajectoire}

\begin{definition}
La trajectoire est l’ensemble des positions successives d’un objet.
\end{definition}

\begin{itemize}
\item rectiligne (ligne droite)
\item circulaire
\item quelconque
\end{itemize}

\begin{example}
\begin{itemize}
\item voiture → rectiligne
\item satellite → circulaire
\item balle → courbe
\end{itemize}
\end{example}

%====================================================
\section{Vitesse}

\begin{definition}
La vitesse moyenne est :
\[
v = \frac{d}{t}
\]
\end{definition}

\begin{example}
Une voiture parcourt 100 km en 2 h :
\[
v = 50 \text{ km/h}
\]
\end{example}

\textbf{Application réelle : radar}

%====================================================
\section{Types de mouvements}

\begin{itemize}
\item uniforme → vitesse constante
\item accéléré → vitesse augmente
\item ralenti → vitesse diminue
\end{itemize}

\begin{example}
\begin{itemize}
\item voiture à vitesse constante → uniforme
\item démarrage → accéléré
\item freinage → ralenti
\end{itemize}
\end{example}

%====================================================
\section{Forces}

\begin{definition}
Une force est une action qui peut modifier le mouvement.
\end{definition}

\textbf{Caractéristiques :}
\begin{itemize}
\item direction
\item sens
\item valeur (en Newton)
\end{itemize}

%====================================================
\section{Le poids}

\begin{definition}
Le poids est la force exercée par la Terre sur un objet.
\end{definition}

\begin{property}
\[
P = m \times g
\]
\end{property}

\begin{example}
Objet de 2 kg :
\[
P = 2 \times 9,8 = 19,6 N
\]
\end{example}

\textbf{Application :}
\begin{itemize}
\item chute d’un objet
\item pesée
\end{itemize}

%====================================================
\section{Forces et équilibre}

\begin{property}
Si les forces se compensent → objet immobile ou mouvement uniforme.
\end{property}

\begin{example}
Livre posé sur une table :
\begin{itemize}
\item poids vers le bas
\item réaction du support vers le haut
\end{itemize}
\end{example}

%====================================================
\section{Forces et mouvement}

\begin{itemize}
\item forces équilibrées → pas de changement
\item forces non équilibrées → mouvement modifié
\end{itemize}

\begin{example}
\begin{itemize}
\item pousser une voiture → accélération
\item freinage → ralentissement
\end{itemize}
\end{example}

%====================================================
\section{Applications concrètes}

\subsection{Sécurité routière}

\begin{itemize}
\item distance de freinage dépend de la vitesse
\item importance des ceintures (forces)
\end{itemize}

\subsection{Sport}

\begin{itemize}
\item trajectoire d’un ballon
\item vitesse d’un coureur
\end{itemize}

\subsection{Espace}

\begin{itemize}
\item mouvement des satellites
\item gravitation
\end{itemize}

%====================================================
\section{Conclusion}

\textbf{À retenir :}
\begin{itemize}
\item mouvement dépend du référentiel
\item vitesse = distance / temps
\item forces expliquent les mouvements
\item poids = force de gravité
\end{itemize}

\bigskip

\begin{center}
\Huge
\textbf{La mécanique permet de comprendre le monde réel}
\end{center}

\end{document}